\scanpage{35}
\unsure{Et|A} cet \deleted{extension provient en fait} élément est en fait ``bien ramifié'' i.e. est dans le sous-groupe
\[
\operatorname{Hom}_{Z_p}(L^2,\underbrace{A^{**}}_{\text{\unsure{Eins-einheiten|1-unités}}})\simeq L^{-2}\otimes_{Z_p}A^{**}.
\]

[\uline{NB.} \unsure{Ça fait donc qu'on a un invariant|On voit donc qu'on a un invariant} à déterminer \uline{mod \deleted{multiplication} \inserted{exponentiation} par un carré de \deleted{\unsure{un|u}} \inserted{unités} $p$-adiques} \dots]

\margin{N.B. L'extension $T_p(E_{\unsure{K|\overline K}})$ peut être triviale (relèvements de Serre) ou non. Cas intéressant: \unsure{si|s'il} est d'ordre fini \dots}
\margin{En tous cas, \unsure{on aura|on a} une opération \uline{abélienne} [dans le cas \unsure{gén.|précédent}]}

b) $E_K$ a bonne réduction, et $E_k$ est ``hypersingulière'' i.e. est non ordinaire. J'ai envie de conjecturer que alors la représentation de $\pi_1(K)$ dans $T_p(E_{\overline K})$ est $\unsure{\frac12|Q}$-simple, voire que l'image de $\pi_1(K)$ dans $\operatorname{Aut}(T_p(E_{\overline K}))$ est un sous-groupe ouvert \dots ??
\margin{\figure{\begin{tikzpicture}\draw (0.7,3) .. controls (0.2,1.5) and (0.2,0.6) .. (0.9,0);\node at (0,1.4) {?};\end{tikzpicture}}}

$(a'+b')$ $E_K$ n'a pas \inserted{\unsure{néc.|néc?}} bonne réduction, mais l'invariant de $E$ est entier, i.e. après extension finie de $K$, on obtient bonne réduction. Ceci splitte en deux cas, a') réduction ordinaire après ext. de la base $K$, i.e. quand on a un invariant ordinaire. Alors l'opération de $\unsure{\pi^1(K)|\pi_1(K)}$ est ``métabélienne''. b') quand on est hypersingulier. Mais \deleted{\unsure{cq|que}} si la conjecture dans b) est valable, les conclusions de cette conjecture valent aussi dans le cas b').
\margin{NB.}
\margin{34}
